% =====================================================================
%  Book 1.  Stimulus-Only Analyses --- Detailed Writeup
%  Auditory Superstitious Perception project
% =====================================================================
\documentclass[11pt]{article}

\usepackage[margin=1in]{geometry}
\usepackage{amsmath,amssymb,amsthm}
\usepackage{graphicx}
\usepackage{booktabs}
\usepackage{caption}
\usepackage{subcaption}
\usepackage{microtype}
\usepackage{xcolor}
\usepackage{listings}
\usepackage[hidelinks]{hyperref}
\usepackage{enumitem}
\usepackage{float}

% ---------- listings setup ----------
\definecolor{codebg}{rgb}{0.96,0.96,0.96}
\definecolor{codekw}{rgb}{0.10,0.30,0.70}
\definecolor{codestr}{rgb}{0.55,0.10,0.10}
\definecolor{codecom}{rgb}{0.30,0.50,0.30}

\lstset{
  language        = Python,
  basicstyle      = \ttfamily\footnotesize,
  backgroundcolor = \color{codebg},
  keywordstyle    = \color{codekw}\bfseries,
  stringstyle     = \color{codestr},
  commentstyle    = \color{codecom}\itshape,
  numbers         = left,
  numberstyle     = \tiny\color{gray},
  numbersep       = 6pt,
  frame           = single,
  framerule       = 0pt,
  breaklines      = true,
  showstringspaces= false,
  columns         = fullflexible,
  tabsize         = 2,
  xleftmargin     = 14pt,
}

\graphicspath{{./}}

% Cap any included figure so it never exceeds the page height.
\newcommand{\figfit}[1]{\includegraphics[width=\linewidth,height=0.75\textheight,keepaspectratio]{#1}}
% Pedagogical convenience boxes
\newcommand{\pintent}[1]{\par\medskip\noindent\textit{Plain-English intent:} #1\par\smallskip}
\newcommand{\pkey}[1]{\par\medskip\noindent\textit{Reader takeaway:} #1\par\smallskip}

\title{\textbf{Book 1.  Stimulus-Only Analyses}\\
       \large What was actually in the auditory superstitious-perception stimulus set?}
\author{Project notebook}
\date{April 2026}

% =====================================================================
\begin{document}
\maketitle
\tableofcontents
\bigskip

% =====================================================================
\section*{Executive summary}
\addcontentsline{toc}{section}{Executive summary}

The stimulus set contains \textbf{300 white-noise clips}, each $L = 2960$
samples long at $8\,\mathrm{kHz}$ ($0.37\,\mathrm{s}$), divided by design into
$150$ \emph{targets} (highly correlated with a fixed template clip
\texttt{wall.wav}) and $150$ \emph{distractors} (uncorrelated with it).  Book~1
asks: \emph{independent of any subject behaviour, what structure does this
set actually contain?}

The short answer:

\begin{itemize}
\item The target/distractor split \emph{is} statistically real under Pearson
      ($d = 0.31$, label-permutation $p < 10^{-3}$), but the effect size is
      tiny: targets are pairwise more similar than distractors by
      $\Delta\bar r \approx 5.8\times 10^{-3}$.
\item That entire effect is the projection onto the template \texttt{wall.wav}.
      After residualising against the template, target--target cohesion
      \emph{vanishes} ($\bar r_{tt} = 3.5\times 10^{-5}$, $p \approx 0.84$).
      The targets do not form an internal family beyond their shared
      wall-component.
\item Under non-Pearson metrics (lagged cross-correlation, envelope
      correlation, log-spectrum correlation, MFCC cosine) the
      target/distractor distinction \emph{collapses} ($d \le 0.04$ for all
      four).  The label is a Pearson-specific artefact of selection.
\item The set has no clean intrinsic clusters: best silhouette across all
      metrics and clustering methods is $0.15$ (envelope, agglomerative),
      and that split isolates $6$ acoustic outliers from the other $294$.
\item Behaviourally, $26$ subjects classified at chance (mean accuracy
      $0.510$, hit rate $0.480$, false-alarm rate $0.460$).  No metric
      predicts trial responses meaningfully (best AUC $= 0.510$).
      Label-blind clusters fit subjects' choices fractionally
      \emph{better} than the experimenter labels do
      ($0.530$ vs.\ $0.510$).
\end{itemize}

The detailed development follows.

% =====================================================================
\section*{How to read this book}
\addcontentsline{toc}{section}{How to read this book}

The four parts walk a single arc:

\begin{itemize}
\item \textbf{Part I (Foundations).} What is on disk: $300$ short white-noise
      clips, half labelled \emph{target} (high Pearson correlation with a
      single fixed template clip \texttt{wall.wav}), half labelled
      \emph{distractor} (near-zero correlation). Everything else in the book
      is built on these centred waveforms.

\item \textbf{Part II (Pearson).} We compare the targets against the
      distractors using the same metric the experimenter used to build the
      set in the first place (Pearson correlation between centred
      waveforms). This is the metric most ``flattering'' to the design,
      and even here the effect is small.

\item \textbf{Part III (other acoustic metrics).} Four other reasonable
      perceptual / acoustic similarity measures (lagged cross-correlation,
      envelope correlation, log-spectrum correlation, MFCC cosine). Under
      these, the target / distractor distinction collapses --- the label is a
      Pearson-specific selection artefact.

\item \textbf{Part IV (synthesis + behavioural sanity check).} A one-pass
      check against subject responses (full Book~2 follow-up is in
      \texttt{book2\_writeup.pdf}) showing that subjects performed at chance
      and that no acoustic metric meaningfully predicts their responses.
\end{itemize}

Each chapter starts with a one-paragraph \emph{Plain-English intent} box, then
gives the formula and the table or plot, and ends with a one-paragraph
\emph{Reader takeaway} box. Skim only the boxes for the narrative; read the
body for derivations and exact numbers.

% =====================================================================
\section{Foundations (Part I)}

\subsection{Stimulus inventory}
\label{sec:inventory}

The stimulus set is organised on disk as
\begin{center}
\ttfamily
raw\_data/audio\_stimuli/\{full\_sentence,imagined\_sentence\}/\{targets,distractors\}/
\end{center}
giving $4$ folders of $75$ \texttt{.wav} files each, for $300$ total.  The
folder split corresponds to the two \emph{blocks} of the experiment:
two folders per block, one for targets and one for distractors.

The selection rule for targets and distractors was Pearson correlation
with a single fixed template, \texttt{raw\_data/audio\_stimuli/wall.wav}.
Concretely, the original pipeline drew $\sim 2{,}000{,}000$ candidate
white-noise clips, computed each candidate's Pearson $r$ to
\texttt{wall.wav}, and selected the top tail as targets and a band near
$r = 0$ as distractors.

\subsection{Preprocessing}

The on-disk files are PCM\_16, $8\,\mathrm{kHz}$, RMS-normalised to
$-23\,\mathrm{dBFS}$ as the final step of the pipeline.  All Book~1
analyses load these on-disk samples directly as \texttt{float64} and
\emph{centre} them (subtract the mean) but do \emph{not} z-score:

\begin{lstlisting}[caption={Canonical loader (excerpt from \texttt{ch11-non-pearson-metrics.py})}]
sr, raw = wavfile.read(wav_path)
rows.append(raw.astype(np.float64))
...
waveform_matrix = np.vstack(rows)
waveform_matrix = waveform_matrix - waveform_matrix.mean(axis = 1, keepdims = True)
\end{lstlisting}

This matters because Pearson correlation is defined on centred signals
but is invariant to scale, and because the energy-based metrics in
Part~III \emph{depend} on absolute amplitude.  Z-scoring would erase
those amplitude differences.

\subsection{Descriptive properties}

Every clip has the same length ($L = 2960$ samples), the same sample
rate ($8000\,\mathrm{Hz}$) and the same nominal RMS.  Mean off-diagonal
Pearson similarity across all $300 \times 300$ pairs is
$+1.40 \times 10^{-3}$ with std $1.86 \times 10^{-2}$:  the clips look
like nearly orthogonal noise vectors in $\mathbb R^{2960}$.

% =====================================================================
\section{Pearson-centered analyses (Part II)}

\subsection{Chapter 4: Template-correlation distribution}
\pintent{Targets were selected by sorting candidate clips on Pearson correlation
with one fixed template, \texttt{wall.wav}. We start by simply looking at the
distribution of those template correlations, both to confirm the selection rule
and to set the scale (in $\mathbb R^{2960}$, what \emph{counts} as a large
correlation?).}

For a centred clip $s \in \mathbb R^L$ and centred template
$t \in \mathbb R^L$, Pearson $r$ is
\begin{equation}
  r(s, t)
  \;=\;
  \frac{\langle s, t \rangle}{\|s\|\,\|t\|}
  \;=\;
  \cos \theta_{st}.
\end{equation}
Selected targets sit in the top tail of $r(\cdot, t_{\text{wall}})$ over
the candidate pool; selected distractors sit near $0$.  In the final set
we have $r(s, t_{\text{wall}}) \in [0.070,\, 0.094]$ for targets and
$|r| \lesssim 10^{-2}$ for distractors (these are the values stored in
\texttt{ch9\_pearson\_r\_to\_template.csv}).

The numerical smallness of even the targets ($r \approx 0.08$, not
$0.5$) is a consequence of the high dimension: a uniformly random unit
vector in $\mathbb R^{2960}$ has $r$ to a fixed unit vector with std
$1/\sqrt{L} \approx 0.0184$, which matches the empirical
distractor-distractor std almost exactly.  Targets sit roughly
$4$--$5\sigma$ out in that null distribution.
\pkey{Targets really are extreme on the template axis ($\sim 4\sigma$), but
the absolute correlations involved are tiny ($r\approx 0.08$). High dimension
makes \emph{small numerical correlations into large statistical
distinctions}; keep that in mind whenever you see $r \sim 10^{-3}$ later.}

\subsection{Chapter 5: Pairwise Pearson structure}
\pintent{If targets share something acoustic, target--target pairs should be
more Pearson-correlated than distractor--distractor or target--distractor
pairs. We compute all three pairwise distributions and compare them.}

Compute the full $300 \times 300$ correlation matrix and split the
upper triangle into three pools by label:
\emph{target--target} ($n=11{,}175$),
\emph{distractor--distractor} ($n=11{,}175$),
\emph{target--distractor} ($n=22{,}500$).

\begin{table}[H]
\centering
\caption{Pairwise Pearson distributions by group.}
\label{tab:pairwise}
\begin{tabular}{lrrrr}
\toprule
Pool & mean & std & vs.\ zero ($t$, $p$) & vs.\ $dd$ (Welch $p$) \\
\midrule
target--target          & $+0.00561$ & $0.0183$ & $32.3,\ 2{\times}10^{-219}$ & $1.7{\times}10^{-119}$ \\
distractor--distractor  & $-0.00017$ & $0.0186$ & $-0.95,\ 0.34$              & --- \\
target--distractor      & $+0.00008$ & $0.0184$ & $0.66,\ 0.51$               & --- \\
\bottomrule
\end{tabular}
\end{table}

Effect sizes: $d_{tt - dd} = 0.313$, $d_{tt - td} = 0.301$,
$d_{dd - td} = -0.013$.  In words: \emph{targets} pairwise correlate
slightly above zero; distractors and cross-pairs are statistically
indistinguishable from zero.

\begin{figure}[H]
\centering
\begin{subfigure}[t]{0.48\textwidth}
  \figfit{stimuli-analyses/outputs/pearson/pariwise/all_data/distributions_histograms.png}
  \caption{Three pairwise distributions overlaid.}
\end{subfigure}\hfill
\begin{subfigure}[t]{0.48\textwidth}
  \figfit{stimuli-analyses/outputs/pearson/pariwise/all_data/correlation_matrix_group_ordered_heatmap.png}
  \caption{$300{\times}300$ correlation matrix, ordered targets-then-distractors.  The faint warm block in the bottom-left is the entire effect.}
\end{subfigure}
\caption{Pairwise Pearson structure (Chapter~5).}
\end{figure}
\pkey{There \emph{is} a target--target excess ($\bar r_{tt} \approx +0.0056$
vs.\ $-0.0002$ for distractors), Cohen's $d \approx 0.31$, and it is
statistically air-tight ($p\sim 10^{-219}$). But the magnitude is the size of
the noise std (both $\sim 0.018$). Targets are not a tight family; they have
a shared mean offset of order $1\%$.}

\subsection{Chapter 6: Target cohesion vs.\ selection rank}
\pintent{If the cohesion in chapter 5 reflects the targets containing more of
the template, the targets that are \emph{most} like the template should be
more cohesive among themselves. Restrict to the top-$k$ most-template-like
targets and recompute mean within-target Pearson, sweeping $k$.}

Order targets by their template correlation $r(s_i, t_{\text{wall}})$,
take the top $k$ for $k \in \{10, 20, \dots, 150\}$, and compute mean
within-pool pairwise Pearson.  If extreme-template targets formed a
tighter family, $\bar r_{tt}^{(k)}$ should rise as $k$ shrinks.

\begin{table}[H]
\centering
\caption{Mean within-target pairwise $r$ for top-$k$ subsets.}
\begin{tabular}{rrr}
\toprule
$k$ & $\bar r_{tt}^{(k)}$ & SE \\
\midrule
$10$  & $0.01189$ & $0.00275$ \\
$20$  & $0.00671$ & $0.00131$ \\
$50$  & $0.00618$ & $0.00051$ \\
$100$ & $0.00565$ & $0.00026$ \\
$150$ & $0.00561$ & $0.00017$ \\
\bottomrule
\end{tabular}
\end{table}

There is a small but real top-tail bump (top~$10$: $\bar r \approx 0.012$,
roughly double the full-set mean), then immediate convergence to the
full-set value by $k \ge 30$.  Selection extremity drives a few extra
percentage points of cohesion in the very top, no more.
\pkey{Top-$10$ targets cohere about twice as strongly as the average target
($\bar r \approx 0.012$ vs.\ $0.006$), but the bump is gone by $k=30$.
Extremely template-aligned targets matter only at the very tail.}

\subsection{Chapter 7: Template alignment vs.\ pairwise cohesion}
\pintent{Now the same idea per-stimulus: does each target's \emph{individual}
template alignment predict how similar it is to the other targets? A simple
scatter and OLS line.}

For each target $i$, plot $r(s_i, t_{\text{wall}})$ against
$\bar r(s_i, S_{\text{tgt}\setminus i})$.  OLS regression gives:
\[
  \text{slope} = 0.0342,\qquad
  R = 0.119,\qquad
  R^2 = 0.014,\qquad
  p = 0.148.
\]
Template alignment explains $\sim\!1.4\%$ of the variation in
how-much-each-target-resembles-the-others.  The relationship is
positive but extremely weak.

\begin{figure}[H]
\centering
\figfit{stimuli-analyses/outputs/pearson/target-cohesion/ch7_template_vs_cohesion.png}
\caption{Per-target template alignment vs.\ mean similarity to other targets.}
\end{figure}
\pkey{The relationship is positive but barely there ($R=0.12$, $R^2=1.4\%$,
$p=0.15$). Template alignment is a small part of why a target resembles other
targets; most variance in per-target cohesion is something else (or noise).}

\subsection{Chapter 8: Pearson-based geometry of the full set}
\pintent{If we take the $300\times 300$ Pearson similarity as a proxy for
stimulus geometry, what does that geometry look like? We embed (MDS, PCA,
$t$-SNE), cluster blindly, and separately ask whether the experimenter's
target--distractor mean-difference direction is a statistically privileged
axis even when no clustering recovers it.}

Build the symmetric distance $D = \mathrm{clip}(1 - R, 0, \cdot)$ from
the full $300{\times}300$ Pearson similarity, and embed.

\paragraph{Classical MDS.}
Double-centre the squared-distance matrix to form the Gram matrix
\[
  B \;=\; -\tfrac{1}{2}\,J\,(D \odot D)\,J,
  \qquad
  J = I - \tfrac{1}{n}\mathbf 1\mathbf 1^\top,
\]
take eigendecomposition $B = U\Lambda U^\top$, and use the top-$k$
positive eigenvectors scaled by $\sqrt{\Lambda_+}$ as coordinates.

\paragraph{Linear PCA on raw waveforms} for comparison.

\paragraph{$t$-SNE} on the precomputed Pearson distance.

\paragraph{Clustering} with agglomerative average-linkage and spectral
$k$-means, $k \in \{2,3,4,5\}$, scored by silhouette and Adjusted Rand
Index (ARI) against experimenter labels.

\paragraph{Label axis.}
Define the unit direction
\[
  \hat u
  \;=\;
  \frac{\bar s_{\text{tgt}} - \bar s_{\text{dist}}}
       {\|\bar s_{\text{tgt}} - \bar s_{\text{dist}}\|},
\]
project every clip onto $\hat u$, and compute Fisher's ratio
$F = (\mu_t - \mu_d)^2 / (\sigma_t^2 + \sigma_d^2)$ and the variance
fraction $\mathrm{Var}(s\cdot\hat u)\,/\,\sum_j \mathrm{Var}(s_{\cdot j})$.
Compare both against $2{,}000$ random unit directions.

\begin{table}[H]
\centering
\caption{Chapter 8 results.}
\begin{tabular}{lrrr}
\toprule
Quantity & Observed & Null mean & $p$ \\
\midrule
MDS top-2 variance fraction        & $1.63\%$ & --- & --- \\
PCA top-2 variance fraction        & $1.11\%$ & --- & --- \\
Best silhouette ($k=2$, spectral)  & $0.0033$ & --- & --- \\
Best ARI vs.\ labels ($k=2$, spectral) & $0.083$ & --- & --- \\
Fisher ratio along label axis      & $37.67$  & $0.009$ & $<10^{-3}$ \\
Variance fraction along label axis & $0.493\%$ & $0.034\%$ & $<10^{-3}$ \\
\bottomrule
\end{tabular}
\end{table}

\textbf{Interpretation.}
The set has no low-dimensional shape: the leading two MDS components
explain only $1.6\%$ of Gram variance.  Generic clustering finds
nothing aligned with labels (best ARI $= 0.08$).  But the
target--distractor mean-difference direction is a statistically
privileged axis ($\sim\!37\times$ stronger than random) that simply
isn't dominant enough ($\sim\!0.5\%$ of waveform variance) to organise
clusters.  Targets are biased in a single shared direction (toward
\texttt{wall.wav}) without being internally tight.

\begin{figure}[H]
\centering
\begin{subfigure}[t]{0.32\textwidth}
\figfit{stimuli-analyses/outputs/pearson/ch8-geometry/ch8_mds.png}
\caption{Classical MDS}
\end{subfigure}
\begin{subfigure}[t]{0.32\textwidth}
\figfit{stimuli-analyses/outputs/pearson/ch8-geometry/ch8_pca.png}
\caption{PCA}
\end{subfigure}
\begin{subfigure}[t]{0.32\textwidth}
\figfit{stimuli-analyses/outputs/pearson/ch8-geometry/ch8_tsne.png}
\caption{$t$-SNE}
\end{subfigure}\\[4pt]
\begin{subfigure}[t]{0.48\textwidth}
\figfit{stimuli-analyses/outputs/pearson/ch8-geometry/ch8_silhouette_ari.png}
\caption{Silhouette and ARI vs.\ $k$}
\end{subfigure}
\begin{subfigure}[t]{0.48\textwidth}
\figfit{stimuli-analyses/outputs/pearson/ch8-geometry/ch8_label_axis_projection.png}
\caption{Projection onto the label axis $\hat u$}
\end{subfigure}
\caption{Chapter 8 outputs.}
\end{figure}
\pkey{Two facts you have to hold in your head simultaneously: (a) generic
blind clustering finds nothing label-related (best ARI $0.08$, best silhouette
$0.003$), but (b) the explicit label-axis direction $\hat u$ is
\emph{statistically privileged} ($F=37.7$, $p<10^{-3}$, $\sim 37\times$
stronger than a random direction). The label is a real but extremely
low-variance direction ($0.49\%$ of waveform variance) --- too weak to drive
clusters but too systematic to be noise.}

\subsection{Chapter 9: Template projection and residual structure}
\pintent{Why does the target--target Pearson cohesion exist at all? Pairwise
Pearson decomposes \emph{exactly} into a contribution from the template
projection $\alpha_i t$ plus a contribution from the residuals $e_i = s_i -
\alpha_i t$. We compute both and see which one carries the cohesion.}

\subsubsection{Decomposition}

Let $t = t_{\text{wall}}$ centred, and $s_i$ centred clip $i$.  The
OLS projection is
\begin{equation}
  \alpha_i \;=\; \frac{\langle s_i, t\rangle}{\langle t, t\rangle},
  \qquad
  \hat s_i = \alpha_i\,t,
  \qquad
  e_i = s_i - \hat s_i.
\end{equation}
Pairwise Pearson decomposes exactly:
\begin{equation}
  \langle s_i, s_j\rangle
  \;=\;
  \alpha_i \alpha_j \langle t, t\rangle
  \;+\;
  \langle e_i, e_j\rangle,
\end{equation}
and dividing by $\sigma_{s_i}\sigma_{s_j}$ gives
\begin{equation}
  r(s_i, s_j)
  \;=\;
  \underbrace{\frac{\alpha_i \alpha_j \,\|t\|^2}{\sigma_{s_i}\sigma_{s_j}}}_{\text{template contribution}}
  \;+\;
  \underbrace{\frac{\langle e_i, e_j\rangle}{\sigma_{s_i}\sigma_{s_j}}}_{\text{residual contribution}}.
\end{equation}

\begin{lstlisting}[caption={Decomposition core (\texttt{ch9-projection-residual.py}).}]
alpha = (X @ t) / (t @ t)              # (n,)
proj  = np.outer(alpha, t)             # (n, L)
resid = X - proj                       # (n, L)
sigma = X.std(axis = 1, ddof = 0)      # (n,)
tmpl_contrib = np.outer(alpha, alpha) * (t @ t) / np.outer(sigma, sigma)
resid_contrib = (resid @ resid.T) / L  / np.outer(sigma, sigma)
# r_total[i,j] == tmpl_contrib[i,j] + resid_contrib[i,j]   (verified <= 1.4e-16)
\end{lstlisting}

\subsubsection{Mean contributions}

\begin{table}[H]
\centering
\caption{Decomposition of mean pairwise Pearson into template and residual parts.}
\begin{tabular}{lrrrr}
\toprule
Pool & $\overline{r}$ & template part & residual part & frac.\ template \\
\midrule
target--target          & $+5.61\!\times\!10^{-3}$ & $+5.58\!\times\!10^{-3}$ & $+3.5\!\times\!10^{-5}$ & $99.4\%$ \\
distractor--distractor  & $-1.67\!\times\!10^{-4}$ & $\sim 0$                 & $-1.67\!\times\!10^{-4}$ & $\sim 0$ \\
target--distractor      & $+8.1\!\times\!10^{-5}$  & $\sim 0$                 & $+8.1\!\times\!10^{-5}$  & $\sim 0$ \\
\bottomrule
\end{tabular}
\end{table}

After residualising the template, target--target cohesion is no longer
distinguishable from zero
($t = 0.20$, $p = 0.84$, Wilcoxon $p = 0.57$, Welch
before-vs-after $p = 2 \times 10^{-112}$).  The targets cohere
\emph{entirely} because they each contain a small piece of
\texttt{wall.wav}; once removed, they revert to noise.

\begin{figure}[H]
\centering
\begin{subfigure}[t]{0.48\textwidth}
\figfit{stimuli-analyses/outputs/pearson/ch9-projection-residual/ch9_pearson_decomposition_bar.png}
\caption{Total $\bar r$ split into template + residual contributions.}
\end{subfigure}\hfill
\begin{subfigure}[t]{0.48\textwidth}
\figfit{stimuli-analyses/outputs/pearson/ch9-projection-residual/ch9_before_after_cohesion.png}
\caption{Before/after residualisation, by group.}
\end{subfigure}
\end{figure}

\subsubsection{Per-sample contributions (Chapter 9b)}

The pairwise contribution at sample $\ell$ is
$c^{(ij)}_\ell = s_{i,\ell}\,s_{j,\ell}$.
Averaging $c^{(ij)}_\ell$ over all target--target pairs gives a per-sample
profile of \emph{which timepoints carry the target--target similarity}.
Correlating that profile against $t_\ell^2$ (template energy),
$|t_\ell|$, and $t_\ell$ (signed):

\begin{table}[H]
\centering
\caption{Per-sample target--target contribution profile vs.\ template features.}
\begin{tabular}{lrr}
\toprule
Comparison & $r$ & $p$ \\
\midrule
$\overline{c^{(tt)}_\ell}$ vs.\ $t_\ell^2$  & $+0.623$ & $\approx 10^{-317}$ \\
$\overline{c^{(tt)}_\ell}$ vs.\ $|t_\ell|$  & $+0.575$ & $\approx 10^{-260}$ \\
$\overline{c^{(tt)}_\ell}$ vs.\ $t_\ell$    & $+0.224$ & $\approx 10^{-35}$ \\
$\overline{c^{(dd)}_\ell}$ vs.\ $t_\ell^2$  & $+0.007$ & $0.70$ \\
\bottomrule
\end{tabular}
\end{table}

The samples that carry target--target similarity are the same samples
where the template is loud.  Distractor pairs show no such pattern.
\pkey{$99.4\%$ of the target--target mean Pearson is the template
contribution $\alpha_i\alpha_j\|t\|^2/(\sigma_i\sigma_j)$. Once that is
subtracted, target--target cohesion is statistically zero ($p=0.84$). Targets
are \emph{not} a coherent acoustic family; they are a set of clips that each
contain a small piece of \texttt{wall.wav}.}

\subsection{Chapter 10: Null models for Pearson structure}
\pintent{The target--target excess is small ($\Delta\bar r \approx 0.006$).
Is even that meaningful? Two complementary nulls: (a) replace all $300$ clips
with fresh white noise and ask whether $\bar r_{tt}$ this big could happen by
chance (white-noise null); (b) keep the data fixed but shuffle the labels and
recompute (label-permutation null).}

Two complementary nulls.

\paragraph{White-noise null.}
Synthesise $300$ clips of length $L$ as iid Gaussian noise, compute
their pairwise Pearson, take mean within-set.  Repeat $500$ times.
\paragraph{Label-permutation null.}
Hold the actual $300{\times}300$ matrix fixed, permute the
target/distractor labels, recompute mean within-tt and within-dd.
Repeat $5{,}000$ times.

\begin{table}[H]
\centering
\caption{Null-model summary (Chapter 10).}
\begin{tabular}{lrrrr}
\toprule
Test & observed & null mean & null std & $p$ \\
\midrule
white-noise null, $\bar r_{tt}$    & $+0.00561$ & $\sim 0$        & $1.7\!\times\!10^{-4}$ & $0.002$ \\
white-noise null, $\bar r_{dd}$    & $-0.00017$ & $\sim 0$        & $1.7\!\times\!10^{-4}$ & $0.33$ \\
label-perm, $\bar r_{tt}$           & $+0.00561$ & $+0.00140$      & $2.2\!\times\!10^{-4}$ & $<2\!\times\!10^{-4}$ \\
label-perm, $\bar r_{tt}-\bar r_{dd}$ & $+0.00578$ & $+6\!\times\!10^{-6}$ & $4.1\!\times\!10^{-4}$ & $<2\!\times\!10^{-4}$ \\
\bottomrule
\end{tabular}
\end{table}

The target--target cohesion sits at the $100$th percentile of both
nulls; distractors are indistinguishable from white noise (which is
exactly what the selection rule tried to achieve).

\begin{figure}[H]
\centering
\begin{subfigure}[t]{0.48\textwidth}
\figfit{stimuli-analyses/outputs/pearson/ch10-null-models/ch10_white_noise_null.png}
\caption{White-noise null.}
\end{subfigure}\hfill
\begin{subfigure}[t]{0.48\textwidth}
\figfit{stimuli-analyses/outputs/pearson/ch10-null-models/ch10_label_permutation_null.png}
\caption{Label-permutation null.}
\end{subfigure}
\end{figure}
\pkey{The observed cohesion sits at the $100$th percentile of both nulls
($p=0.002$ white-noise, $p<2\times 10^{-4}$ label-permutation). The
distractor pool is indistinguishable from white noise, exactly as the
selection rule intended.}

% =====================================================================
\section{Non-Pearson metric analyses (Part III)}

\subsection{Chapter 11: Choosing alternative metrics}
\pintent{The selection rule was Pearson on raw centred waveforms, but a human
listener perceives audio through filters that look very little like that.
We define four perceptually-motivated similarity measures (envelope-based,
spectral, cepstral, time-shift-tolerant) so that subsequent chapters can ask
the label question through several different lenses.}

Five similarity measures, each chosen for an interpretable axis of
acoustic structure.

\paragraph{1. Pearson raw} (reference).
$r(s_i, s_j) = \cos(s_i, s_j)$ on centred waveforms.

\paragraph{2. Lagged cross-correlation, $\pm 200\,\mathrm{ms}$.}
Full sample-domain cross-correlation via FFT:
$x_{ij}(\tau) = \sum_t s_{i,t} s_{j,t+\tau}\,/\,(\|s_i\|\|s_j\|)$,
report $\max_{|\tau|\le \tau_{\max}} |x_{ij}(\tau)|$.  Captures similarity
that survives small temporal shifts.

\paragraph{3. Envelope correlation.}
Compute the analytic-signal magnitude (Hilbert envelope), low-pass with
a $161$-sample moving average ($\sim 50\,\mathrm{Hz}$ cutoff at $8\,\mathrm{kHz}$),
then Pearson on envelopes.  Removes carrier phase, keeps loudness
contour.

\paragraph{4. Log-spectrum correlation.}
$\tilde s = \log(1 + |\mathcal F(s)|)$ from the real FFT, then Pearson
across frequency bins.  Compares spectral shape, ignores phase.

\paragraph{5. MFCC cosine.}
$12$ Mel-frequency cepstral coefficients per frame (dropping $C_0$,
which collapses to a constant under RMS normalisation), summarised by
mean and std across frames, $z$-scored per dimension across clips,
cosine similarity.  Approximates a coarse speech-style timbral
descriptor.

\begin{lstlisting}[caption={Lagged cross-correlation core.}]
F = next_fast_len(2*L - 1)
S = np.fft.rfft(X, n = F, axis = 1)            # (n, F/2+1)
norms = np.linalg.norm(X, axis = 1)
for i in range(n):
    cc = np.fft.irfft(S[i] * np.conj(S), n = F, axis = 1)
    cc = np.concatenate([cc[:, -max_lag:], cc[:, :max_lag+1]], axis = 1)
    out[i, :] = np.max(np.abs(cc), axis = 1) / (norms[i] * norms)
\end{lstlisting}
\pkey{Five orthogonal lenses on the same $300$ clips. Pearson is the
selection metric; the other four ask whether the label structure shows up
under time-shift-tolerant, envelope, spectral, or cepstral similarity.}

\subsection{Chapter 12: Label-aware separation under each metric}
\pintent{Repeat the chapter 5 target--target / distractor--distractor split
for each non-Pearson metric. If any of these acoustic descriptions agreed
with the label, we would see a non-trivial Cohen's $d$.}

Same target--target / distractor--distractor / target--distractor
decomposition as Chapter~5, repeated for each metric.

\begin{table}[H]
\centering
\caption{Cohen's $d$ for $tt$ vs.\ $dd$ under each metric.}
\label{tab:ch12}
\begin{tabular}{lrrr}
\toprule
Metric & $\bar r_{tt}$ & $\bar r_{dd}$ & $d_{tt-dd}$ \\
\midrule
pearson\_raw       & $+0.00561$ & $-0.00017$ & $\mathbf{0.313}$ \\
lagged\_xcorr\_max & $+0.06186$ & $+0.06165$ & $0.037$ \\
envelope\_corr     & $+0.74480$ & $+0.74344$ & $0.022$ \\
log\_spectrum\_corr & $+0.17737$ & $+0.17726$ & $0.005$ \\
mfcc\_cosine       & $-0.00072$ & $-0.00137$ & $0.003$ \\
\bottomrule
\end{tabular}
\end{table}

\textbf{Interpretation.}
Pearson is the only metric under which the target/distractor split is
visible at all.  Under envelope correlation the average pair has
$r \approx 0.74$ regardless of label --- all clips look alike to the
envelope, because they all share the same RMS-normalised amplitude
contour profile.  The label is a property of the metric, not of the
stimuli.

\begin{figure}[H]
\centering
\figfit{stimuli-analyses/outputs/non_pearson/ch12-label-aware/ch12_relative_to_pearson.png}
\caption{Per-metric label separation, normalised against Pearson.}
\end{figure}
\pkey{Pearson is the only metric under which the label is even faintly
visible ($d=0.31$). Lagged xcorr, envelope, log-spectrum, MFCC all give
$d \le 0.04$. The target/distractor distinction is a Pearson artefact, not
an acoustic property a listener could realistically pick up.}

\subsection{Chapter 13: Label-blind clustering under each metric}
\pintent{Forget the labels and ask: under each metric, does the set break
into natural clusters at $k=2$? If acoustic structure exists that is not the
experimenter's split, blind clustering should find it.}

Convert similarity $s_{ij} \in [-1,1]$ to distance $d_{ij} = \mathrm{clip}(1 - s_{ij}, 0, 2)$
and to affinity $a_{ij} = (s_{ij} + 1)/2$.  Cluster at $k=2$ with
\begin{itemize}
\item agglomerative average linkage on $D$,
\item spectral $k$-means on $A$,
\end{itemize}
and score with silhouette (intrinsic) and ARI vs.\ experimenter labels.

\begin{table}[H]
\centering
\caption{Chapter 13 results, $k=2$.}
\begin{tabular}{llrrr}
\toprule
Metric & Method & sizes & silhouette & ARI \\
\midrule
pearson\_raw       & agglomerative & $[287, 13]$ & $0.003$ & $0.005$ \\
pearson\_raw       & spectral      & $[156, 144]$ & $0.003$ & $0.083$ \\
lagged\_xcorr\_max & agglomerative & $[256, 44]$  & $0.001$ & $-0.001$ \\
lagged\_xcorr\_max & spectral      & $[166, 134]$ & $0.001$ & $0.022$ \\
envelope\_corr     & agglomerative & $[294, 6]$   & $\mathbf{0.151}$ & $-0.000$ \\
envelope\_corr     & spectral      & $[118, 182]$ & $0.087$ & $-0.003$ \\
log\_spectrum\_corr& agglomerative & $[298, 2]$   & $0.011$ & $-0.000$ \\
log\_spectrum\_corr& spectral      & $[145, 155]$ & $0.005$ & $0.000$ \\
mfcc\_cosine       & agglomerative & $[203, 97]$  & $0.055$ & $-0.001$ \\
mfcc\_cosine       & spectral      & $[154, 146]$ & $0.086$ & $-0.003$ \\
\bottomrule
\end{tabular}
\end{table}

The best silhouette ($0.15$, envelope agglomerative) corresponds to a
$294/6$ split, i.e.\ peeling six acoustic outliers off the rest.  No
clustering recovers the experimenter labels (max ARI $0.083$).
\pkey{The best cluster split (envelope agglomerative, silhouette $0.15$)
peels off six acoustic outliers from the other $294$. There is no clean
$k=2$ partition under any metric, and certainly none aligned with the
experimenter's labels.}

\subsection{Chapter 14: Cross-metric synthesis}
\pintent{Bring the five metrics onto a single scoreboard along three axes
(target cohesion, label separation, intrinsic clustering), and check whether
the metrics describe redundant geometries (Mantel-style cross-metric
correlations).}

Three axes summarise each metric:
\begin{enumerate}[leftmargin=*]
\item \textbf{Target cohesion} (z-score of $\bar r_{tt}$ against the
metric-specific off-diagonal distribution),
\item \textbf{Label separation} (Cohen's $d_{tt - dd}$),
\item \textbf{Intrinsic clustering} (best silhouette over methods at $k=2$).
\end{enumerate}

\begin{table}[H]
\centering
\caption{Cross-metric ranking.}
\begin{tabular}{lrrr}
\toprule
Metric & cohesion ($z$) & label sep.\ ($d$) & intrinsic silh. \\
\midrule
pearson\_raw         & $\mathbf{0.227}$ & $\mathbf{0.313}$ & $0.003$ \\
lagged\_xcorr\_max   & $0.017$ & $0.037$ & $0.001$ \\
envelope\_corr       & $0.014$ & $0.022$ & $\mathbf{0.151}$ \\
log\_spectrum\_corr  & $0.006$ & $0.005$ & $0.011$ \\
mfcc\_cosine         & $0.012$ & $0.003$ & $0.086$ \\
\bottomrule
\end{tabular}
\end{table}

\paragraph{Mantel-style cross-metric similarity.}
Flatten each metric's upper-triangular pair vector and compute
Pearson's $r$ between metrics.  Maximum off-diagonal $|r| = 0.078$
(lagged vs.\ log-spectrum), most $|r| < 0.01$.  The five metrics are
nearly orthogonal as ways of describing this set.  No metric is a
disguised duplicate of another.

\begin{figure}[H]
\centering
\begin{subfigure}[t]{0.55\textwidth}
\figfit{stimuli-analyses/outputs/non_pearson/ch14-comparison/ch14_axis_rankings.png}
\caption{Per-metric rankings on the three axes.}
\end{subfigure}\hfill
\begin{subfigure}[t]{0.40\textwidth}
\figfit{stimuli-analyses/outputs/non_pearson/ch14-comparison/ch14_cross_metric_mantel.png}
\caption{Cross-metric Mantel correlations.}
\end{subfigure}
\end{figure}
\pkey{The five metrics are nearly orthogonal as descriptions of this set
(max cross-metric $|r|=0.078$). No metric is a redundant copy of another;
each really is asking a different question. The label is visible only to
Pearson; intrinsic clustering structure is visible only to envelope.}

% =====================================================================
\section{Synthesis and behavioural cross-check (Part IV)}

\subsection{Chapter 15: What did subjects do?}
\pintent{Even though Book~1 is about the stimulus set, we owe ourselves a
one-pass behavioural cross-check: do any of these metrics relate to subject
responses? If not, that is itself the most informative possible result.
Full subject-side analysis lives in \texttt{book2\_writeup.pdf}.}

The Book~1 mandate is stimulus-only, but a one-pass behavioural
sanity-check answers the question \emph{``does any of this matter for
how subjects responded?''} without committing to a full Book~2.

\paragraph{Trials.}
$26$ subjects across two sites, $7{,}748$ trials total
($\sim 2 \times 150$ trials per subject).  Each trial yields
\texttt{(stimulus\_number, true\_label, subject\_response)}.

\paragraph{Per-stimulus feature.}
For each metric $M$ and each clip $i$ in block $B$,
\[
  \Delta^M_i
  \;=\;
  \overline{M(s_i, s_j)}_{j\in B,\, \text{label}_j = \text{tgt},\, j \ne i}
  \;-\;
  \overline{M(s_i, s_j)}_{j\in B,\, \text{label}_j = \text{dist},\, j \ne i}.
\]
A clip with large positive $\Delta^M$ is more like targets than
distractors under $M$.

\paragraph{Stimulus-level prediction.}
Pearson $r$ between $\Delta^M_i$ and the empirical
$p_i = \Pr(\text{subject called } s_i \text{ a target})$:

\begin{table}[H]
\centering
\begin{tabular}{lrr}
\toprule
Metric & $r$ & $p$ \\
\midrule
pearson\_raw        & $+0.094$ & $0.11$ \\
envelope\_corr      & $+0.078$ & $0.18$ \\
mfcc\_cosine        & $+0.059$ & $0.31$ \\
log\_spectrum\_corr & $+0.044$ & $0.45$ \\
lagged\_xcorr\_max  & $-0.007$ & $0.91$ \\
\bottomrule
\end{tabular}
\end{table}

\paragraph{Trial-level prediction (logistic AUC).}
Predict $\Pr(\text{called target})$ from $\Delta^M$ alone, vs.\ from
the experimenter's true label:

\begin{table}[H]
\centering
\begin{tabular}{lr}
\toprule
Predictor & AUC \\
\midrule
true\_label\_baseline & $0.510$ \\
envelope\_corr        & $0.510$ \\
mfcc\_cosine          & $0.507$ \\
log\_spectrum\_corr   & $0.504$ \\
lagged\_xcorr\_max    & $0.502$ \\
pearson\_raw          & $0.490$ \\
\bottomrule
\end{tabular}
\end{table}

\paragraph{Cluster vs.\ response agreement.}
Match each metric's label-blind $k=2$ clustering against the binary
subject response, picking the better of the two polarities:

\begin{table}[H]
\centering
\begin{tabular}{llr}
\toprule
Metric & method & match rate \\
\midrule
log\_spectrum\_corr & agglomerative & $\mathbf{0.530}$ \\
envelope\_corr      & agglomerative & $0.529$ \\
pearson\_raw        & agglomerative & $0.526$ \\
mfcc\_cosine        & agglomerative & $0.524$ \\
lagged\_xcorr\_max  & agglomerative & $0.522$ \\
true\_label\_baseline & ---         & $0.510$ \\
\bottomrule
\end{tabular}
\end{table}

\begin{figure}[H]
\centering
\begin{subfigure}[t]{0.48\textwidth}
\figfit{stimuli-analyses/outputs/pearson/ch15-subject-prediction/ch15_metric_trial_auc.png}
\caption{Trial-level AUC.}
\end{subfigure}\hfill
\begin{subfigure}[t]{0.48\textwidth}
\figfit{stimuli-analyses/outputs/pearson/ch15-subject-prediction/ch15_response_agreement.png}
\caption{Cluster-vs-response agreement.}
\end{subfigure}
\end{figure}

\textbf{Interpretation.}
\begin{enumerate}[leftmargin=*]
\item Subjects performed at chance.  Mean accuracy $0.510$, hit rate
$0.480$, false-alarm rate $0.460$.
\item No metric, including the experimenter labels, predicts trial
responses meaningfully.
\item The label-blind agglomerative clusters consistently beat the
experimenter labels by $\sim 2$ percentage points.  These clusters tend
to isolate small acoustic-outlier groups, suggesting subjects were
weakly tracking gross acoustic salience rather than the imposed
identity contrast.
\end{enumerate}
\pkey{Subjects were at chance (mean accuracy $0.510$). No metric --- not even
the experimenter labels --- predicts trial responses meaningfully (best AUC
$0.510$). Surprisingly, blind acoustic clusters fit subject responses
\emph{better} than the experimenter labels, by about $2$ points. This
motivated Book~2.}

% =====================================================================
\section{What the stimulus set actually contains}

Pulling it all together:

\begin{description}
\item[Structure imposed by design.]
A small ($\Delta\bar r \approx 0.006$, $d \approx 0.31$) but
statistically robust target--distractor separation under Pearson, driven
\emph{entirely} by shared projection onto \texttt{wall.wav}.
\item[Structure actually present.]
Beyond the wall-projection direction, essentially noise.  Targets are
not internally cohesive once the template is removed.  No metric
reveals interesting intrinsic clusters; the set is a $300$-point cloud
in a high-dimensional space with one weak preferred direction.
\item[Structure visible only under some metrics.]
Only Pearson sees the label.  Under any other reasonable acoustic
similarity (envelope, lagged $r$, log-spectrum, MFCC), targets and
distractors are statistically indistinguishable.  Subjects, who
presumably do not perform sample-level Pearson against the template in
their heads, would have no acoustic basis to recover the labels.
\item[Implications for subject analyses.]
Behavioural performance was at chance, consistent with the conclusion
above.  The interesting variance for downstream books is between
\emph{subjects}, not between stimuli.  Per-stimulus prediction of
behaviour by acoustic features is a dead end with this stimulus set.
\end{description}

% =====================================================================
\appendix
\section{File map}

\paragraph{Scripts (all under \texttt{stimuli-analyses/}).}
\begin{itemize}[leftmargin=*]
\item \texttt{pairwise-pearson.py} --- Chapter~5
\item \texttt{target-cohesion-via-rank.py} --- Chapters 6, 7
\item \texttt{ch8-pearson-geometry.py} --- Chapter 8
\item \texttt{ch9-projection-residual.py} --- Chapter 9
\item \texttt{ch9b-per-sample-contributions.py} --- Chapter 9b
\item \texttt{ch10-null-models.py} --- Chapter 10
\item \texttt{ch11-non-pearson-metrics.py} --- Chapter 11 (caches metrics)
\item \texttt{ch12-label-aware-non-pearson.py} --- Chapter 12
\item \texttt{ch13-label-blind-non-pearson.py} --- Chapter 13
\item \texttt{ch14-metric-comparison.py} --- Chapter 14
\item \texttt{ch15-subject-response-prediction.py} --- Chapter 15
\end{itemize}

\paragraph{Outputs.}
Tables and figures live under
\texttt{stimuli-analyses/outputs/\{pearson,non\_pearson\}/<chapter>/}.

\section{Key formulae cheat sheet}

\begin{align*}
r(s, t)
  &\;=\; \frac{\langle s_c, t_c\rangle}{\|s_c\|\,\|t_c\|}
   \qquad (\text{centred Pearson}) \\
\alpha_i
  &\;=\; \frac{\langle s_i, t\rangle}{\langle t, t\rangle}
   \qquad (\text{OLS template projection}) \\
e_i
  &\;=\; s_i - \alpha_i t
   \qquad (\text{residual}) \\
r(s_i, s_j)
  &\;=\; \frac{\alpha_i\alpha_j\|t\|^2 + \langle e_i, e_j\rangle}{\sigma_{s_i}\sigma_{s_j}} \\
F
  &\;=\; \frac{(\mu_t - \mu_d)^2}{\sigma_t^2 + \sigma_d^2}
   \qquad (\text{Fisher discriminant ratio}) \\
\hat u
  &\;=\; \frac{\bar s_{\text{tgt}} - \bar s_{\text{dist}}}{\|\bar s_{\text{tgt}} - \bar s_{\text{dist}}\|}
   \qquad (\text{label axis}) \\
B
  &\;=\; -\tfrac{1}{2}\,J\,(D \odot D)\,J,\quad J = I - n^{-1}\mathbf 1\mathbf 1^\top
   \qquad (\text{classical MDS Gram}) \\
d
  &\;=\; \frac{\mu_a - \mu_b}{\sqrt{\tfrac{1}{2}(\sigma_a^2 + \sigma_b^2)}}
   \qquad (\text{Cohen's }d).
\end{align*}

\end{document}
